% 4. 惯性 Langlands 对应（底本 p.32--35）
\section{惯性 Langlands 对应}

\subsection{Bernstein 分解}

\subsubsection{Bernstein 分量与分解}

\begin{definition}
设 $L\subset\GL_n(K)$ 是 Levi 子群，$P\subset\GL_n(K)$ 是以 $L$ 为 Levi 因子的抛物子群。
设 $\varsigma$ 是 $L$ 的不可约超尖点表示。\emph{Bernstein 分量}
\[
R_{(L,\varsigma)}=R_{(L,\varsigma)}\bigl(\GL_n(K)\bigr)\subset\Repinf_n
\]
是这样一个满子范畴：$\pi$ 属于它，当且仅当 $\pi$ 的每个不可约子商都是
\[
I^{\GL_n}_L(\varsigma\otimes\eta):=\Ind^{\GL_n(K)}_P(\varsigma\otimes\eta)
\]
的子商，其中 $\eta$ 是 $L$ 的非分歧特征标，而 $I^{\GL_n}_L:=\Ind^{\GL_n(K)}_P\circ j^{*}$
是正规化抛物诱导，$j\colon P\to L$ 是到 Levi 商的映射。
\end{definition}

\begin{lemma}
两个 Bernstein 分量 $R_{(L_1,\varsigma_1)}$ 与 $R_{(L_2,\varsigma_2)}$ 相等，
当且仅当它们\emph{惯性等价}，即存在 $L_2$ 的非分歧特征标 $\eta$ 与 $g\in\GL_n(K)$ 使
\[
gL_1g^{-1}=L_2,\qquad g\varsigma_1=\varsigma_2\eta .
\]
\end{lemma}

\begin{definition}
定义 $\GL_n(K)$ 的\emph{Bernstein 中心}为 Bernstein 分量的集合，即
\[
\mathcal{B}=\mathcal{B}\bigl(\GL_n(K)\bigr):=\{(L,\varsigma)\}/\text{惯性等价}.
\]
\end{definition}

\begin{theorem}[Bernstein 分解]
$\GL_n(K)$ 的光滑表示范畴 $\Repinf_n$ 分解为直积
\[
\Repinf_n=\prod_{(L,\varsigma)\in\mathcal{B}}R_{(L,\varsigma)} .
\]
\end{theorem}

\begin{remark}
要确定一个给定的不可约表示 $\pi\in\An(K)$ 落在哪个 Bernstein 分量中，
只需考虑它的超尖点支撑。
\end{remark}

\subsubsection{$\GL_2$ 的 Bernstein 分量}

\begin{example}[$n=2$]
Bernstein 分量 $R_{(L,\sigma)}(\GL_2(K))$ 是下列类型之一：
\begin{enumerate}[label=(\arabic*),leftmargin=2.2em]
\item 若 $R_{(L,\sigma)}(G)$ 含有不可约超尖点表示 $\pi$，则 $(L,\sigma)$ 与 $(G,\pi)$
      惯性等价，而 Bernstein 分量 $R_{(G,\pi)}(G)$ 中的不可约表示为
\[
\{\pi\otimes(\eta\circ\det):\ \eta\ \text{非分歧}\}.
\]
\item 若 $R_{(L,\sigma)}(G)$ 含有 $\pi=\St(\chi)=\St_G\otimes(\chi\circ\det)$，
      则 $(L,\sigma)\sim(T,\chi\otimes\chi)$，而 $R_{(T,\chi)}(G)$ 中的不可约表示为
\[
\bigl\{\St_G\otimes(\chi\eta\circ\det)\bigr\}
\ \cup\
\bigl\{(\eta_1\chi|\cdot|^{-1/2})\times(\eta_2\chi|\cdot|^{1/2})\bigr\},
\]
其中 $\eta,\eta_1,\eta_2$ 取遍非分歧特征标（且使后一族仍为主序列）。
\item 若 $\pi=\pi(\chi_1,\chi_2)$，其中 $\chi_1\chi_2^{-1}$ 分歧
      （否则落入前一情形），则该 Bernstein 分量中的不可约表示为
\[
\{(\eta_1\chi_1)\times(\eta_2\chi_2)\},
\]
其中 $\eta_1,\eta_2$ 非分歧。
\end{enumerate}
\end{example}

\begin{remark}
在差一个非分歧特征标的意义下，上例中各个 Bernstein 分量对应的超尖点支撑为：
\begin{enumerate}[label=(\arabic*),leftmargin=2.2em]
\item $R_{(G,\pi)}(G)$ 对应 $\{\pi\}$；
\item $R_{(T,\chi\otimes\chi)}(G)$ 对应 $\{\chi,\chi\}$；
\item $R_{(T,\chi_1\otimes\chi_2)}(G)$ 对应 $\{\chi_1,\chi_2\}$，其中 $\chi_1\chi_2^{-1}$ 分歧。
\end{enumerate}
\end{remark}

\subsubsection{光滑表示的超尖点部分}

底本此处只有标题，正文未写。

\subsubsection{Bernstein 分解的证明}

底本此处只有标题，正文未写。

\subsection{惯性对应}

\subsubsection{Weil--Deligne 惯性类型}

先需要一个（有限维）表示论中的标准引理：

\begin{lemma}
设 $H\triangleleft G$ 是正规子群，$G/H$ 交换，$(\rho_1,V_1)$ 与 $(\rho_2,V_2)$ 是
$\mathbb{C}$ 上（或任一代数闭域上）两个不可约有限维表示，且 $\rho_1|_H\cong\rho_2|_H$，
则存在特征标 $\chi\colon G\to G/H\to\mathbb{C}^{\times}$，使 $\rho_2\cong\rho_1\otimes\chi$。
\end{lemma}

\begin{proof}[略证]
把 $\Hom_{\mathbb{C}}(V_1,V_2)$ 看作带有 $G$-作用的向量空间。由
$\rho_1|_H\cong\rho_2|_H$ 知
\[
W:=\Hom_H(\rho_1,\rho_2)=\bigl(\Hom_{\mathbb{C}}(V_1,V_2)\bigr)^{H}\neq 0 .
\]
$W\subset\Hom_{\mathbb{C}}(V_1,V_2)$ 是 $G$ 的子表示，且 $H$ 在其上平凡作用，
于是交换商 $G/H$ 作用在 $W$ 上；$W$ 因此含有一个一维 $G$-子表示：
\[
\chi\hookrightarrow W\subset V_1^{\vee}\otimes V_2 .
\]
特别地，用 $\chi$ 扭转可得
\[
\triv\hookrightarrow(V_1\otimes\chi)^{\vee}\otimes V_2,
\]
即
\[
\triv\hookrightarrow\bigl((V_1\otimes\chi)^{\vee}\otimes V_2\bigr)^{G}
=\Hom_G(V_1\otimes\chi,V_2)\neq 0 .
\]
于是由 Schur 引理，$\rho_1\chi\cong\rho_2$。
\end{proof}

\begin{remark}
事实上我们证明了稍强一些的断言，下文会用到：在上述设定中，若 $\rho_1$ 与 $\rho_2$
都不可约，且 $\Hom_H(\rho_1,\rho_2)\neq 0$，则同样的结论成立。
也就是说，引理中并不需要事先假设 $\rho_1|_H\cong\rho_2|_H$。
\end{remark}

\begin{lemma}
$G$ 的两个不可约容许表示 $\pi$ 与 $\pi'$ 属于同一个 Bernstein 分量
$R_{(L,\sigma)}(G)$，当且仅当由 LLC 得到的（Weil--Deligne 表示的）底层 $W_K$-表示满足
\[
\rec_n(\pi)|_{I_K}\cong\rec_n(\pi')|_{I_K}.
\]
换言之，Bernstein 中心 $\mathcal{B}(\GL_n)$ 由 Weil--Deligne 惯性类型标定。
\end{lemma}

\begin{proof}
由 Bernstein--Zelevinsky 分类，$\pi=Q(\Delta_1,\dots,\Delta_r)$，
其中 $\Delta_i=\Delta_i(\pi_i,m_i)$ 是 Zelevinsky 段，于是 $\pi$ 属于这样的 Bernstein 分量：
\[
L=\prod_{i=1}^{r}\GL_{n_i}(K)^{m_i},
\qquad
\sigma=\bigotimes_{i=1}^{r}\ \bigotimes_{j=0}^{m_i-1}\pi_{i,j},
\]
其中 $\pi_{i,j}$ 是 $\GL_{n_i}(K)$ 的超尖点表示 $\pi_i|\cdot|^{j}$。由局部 Langlands 对应，
\[
\rec_n(\pi)=\bigoplus_{i=1}^{r}\Sp_{m_i}\bigl(\rec_{n_i}(\pi_i)\bigr),
\]
而限制到 $I_K$ 的底层 $W_K$-表示为
\[
\rec_n(\pi)|_{I_K}=\bigoplus_{i=1}^{r}\rec_{n_i}(\pi_i)_{m_i}|_{I_K}.
\]
于是归结到超尖点情形：$\GL_k(K)$ 的两个超尖点表示 $\pi$ 与 $\pi'$ 属于同一
Bernstein 分量（即二者相差一个非分歧特征标 $\eta$ 的扭转），当且仅当
$\rec(\pi)|_{I_K}\cong\rec(\pi')|_{I_K}$，而这可由引理 4.3 推出。
注意最后一步隐含地用到了引理 4.3 之后的注：设
\[
\bigoplus_{i=1}^{r}\rec_{n_i}(\pi_i)_{m_i}|_{I_K}\cong
\bigoplus_{i=1}^{r'}\rec_{n'_i}(\pi'_i)_{m'_i}|_{I_K},
\]
我们（逐次）应用该注，可得对某些非分歧特征标 $\psi_i$ 有
$\rec_{n_i}(\pi_i)\otimes\psi_i\cong\rec_{n_i}(\pi'_i)$。
\end{proof}

\subsubsection{Galois 惯性类型与一般表示}

设 $\mathcal{K}=\GL_n(\mathcal{O}_K)$ 为 $\GL_n(K)$ 的标准极大紧开子群。
（底本此处仍写作 $K=\GL_n(\mathcal{O}_K)$，与域 $K$ 撞记号；译文改用 $\mathcal{K}$。）

\begin{lemma}
设 $\pi,\pi'$ 是 $\GL_n(K)$ 的两个一般不可约容许表示，其相应的 $l$-进表示满足
\[
r_n(\pi)|_{I_K}\cong r_n(\pi')|_{I_K},
\qquad r_n=\mathrm{WD}^{-1}\circ\rec_n .
\]
则
\[
\pi|_{\mathcal{K}}\cong\pi'|_{\mathcal{K}} .
\]
换言之，若两个一般表示有相同的 Galois 惯性类型，则它们限制到 $\mathcal{K}$ 上同构。
\end{lemma}

\begin{proof}
由 Iwasawa 分解 $\GL_n(K)=B\cdot\GL_n(\mathcal{O}_K)$（$B$ 为上三角矩阵构成的 Borel 子群）
以及诱导表示限制的 Mackey 分解，对任一（标准）抛物子群 $P\subset\GL_n(K)$
与 $P$ 的表示 $\sigma$，有
\[
\bigl(\Ind^{\GL_n(K)}_P\sigma\bigr)|_{\mathcal{K}}
=\Ind^{\mathcal{K}}_{P\cap\mathcal{K}}\bigl(\sigma|_{P\cap\mathcal{K}}\bigr).
\]
因 $\pi$ 一般，由定理 3.1 有 $\pi=Q(\Delta_1)\times\cdots\times Q(\Delta_r)$，且
$\rec_n(\pi)=\bigoplus_i\rec_{n_im_i}(Q(\Delta_i))$，$\pi'$ 同理。$W_K$ 的两个表示
\[
\bigoplus_{i=1}^{r}\Sp_{m_i}(\rho_i),
\qquad
\bigoplus_{i=1}^{r'}\Sp_{m'_i}(\rho'_i)
\]
有相同的 Galois 惯性类型，当且仅当由引理 1.9，$r=r'$，且在重排之后
$m_i=m'_i$、$\rho_i|_{I_K}\cong\rho'_i|_{I_K}$。

于是只需在 $\pi=Q(\Delta(\lambda,m))$ 与 $\pi'=Q(\Delta(\lambda',m))$ 的假设下证明本引理，
其中 $\lambda,\lambda'$ 是 $\GL_d(K)$ 的超尖点表示。注意 $\lambda=\lambda'\otimes\eta$，
其中 $\eta$ 是 $\GL_d(K)$ 的非分歧特征标。特征标 $\eta$ 形如 $\eta=\chi\circ\det$，
$\chi$ 是 $K^{\times}\to\mathbb{C}^{\times}$ 的非分歧拟特征标。于是
$\rec_n(\pi)=\rec_n(\pi')\otimes\rec_1(\chi)$，由此得 $\pi=\pi'\chi$。
\end{proof}

\begin{remark}
「一般」这一条件不可省。例如对 $\GL_2(K)$，$\triv$ 与 $\St_2$ 落在同一
Bernstein 分量中，但它们限制到 $\mathcal{K}=\GL_2(\mathcal{O}_K)$ 上并不同构。
\end{remark}

\subsection{类型}

最后我们简述类型论。

\subsubsection{标准定义}

\begin{definition}
设 $\sigma$ 是 $\mathcal{K}$ 的光滑不可约表示，$s=(L,\varsigma)\in\mathcal{B}(\GL_n)$
是 Bernstein 中心中的一个元素。
\begin{enumerate}[label=(\arabic*),leftmargin=2.2em]
\item 若对一切光滑不可约表示 $\pi\in\Repinf_n$，
\[
\pi|_{\mathcal{K}}\ \text{包含}\ \sigma
\quad\Longleftrightarrow\quad
\pi\in R_{(L,\varsigma)},
\]
则称 $\sigma$ 是 $(L,\varsigma)$ 的\emph{类型}（type）。
\item 若对一切光滑不可约表示 $\pi$，
\[
\pi|_{\mathcal{K}}\ \text{包含}\ \sigma
\quad\Longrightarrow\quad
\pi\in R_{(L,\varsigma)},
\]
则称 $\sigma$ 是 $(L,\varsigma)$ 的 \emph{$\mathcal{K}$-型}。
\item 设 $\pi\in R_{(L,\varsigma)}$ 是 $R_{(L,\varsigma)}$ 中的光滑不可约表示，
      若 $\sigma$ 是 $(L,\varsigma)$ 的 $\mathcal{K}$-型，且在 $\pi|_{\mathcal{K}}$ 中出现的
      $\mathcal{K}$-型里维数最小，则称 $\sigma$ 是 $\pi$ 的\emph{极小 $\mathcal{K}$-型}。
\end{enumerate}
\end{definition}

注意每个类型显然都是 $\mathcal{K}$-型，且这两个概念只依赖于
$(L,\varsigma)\in\mathcal{B}(\GL_n)$；而 $\pi$ 的极小 $\mathcal{K}$-型的定义依赖于
特定的 $\pi\in R_{(L,\varsigma)}$。也就是说，对给定的 $s=(L,\varsigma)$ 可以有很多
$\mathcal{K}$-型；对每个不可约 $\pi\in R_{(L,\varsigma)}$，一个 $\mathcal{K}$-型 $\sigma$
可能在 $\pi|_{\mathcal{K}}$ 中出现、也可能不出现；在那些出现的当中，维数最小者
就是 $\pi$ 的极小 $\mathcal{K}$-型。

\begin{remark}
我们对「类型」的定义与 [9] 中的定义一致，而与 [4] 中的定义略有不同。
后者我们称为 \emph{Bushnell--Kutzko 型}，它是一个对 $(J,\sigma)$，
其中 $J$ 是紧开子群、$\sigma$ 是 $J$ 的光滑不可约表示——
也就是说，它允许比 $\mathcal{K}=\GL_n(\mathcal{O}_K)$ 更一般的紧开子群。
\end{remark}

\begin{keytakeaways}
\begin{itemize}[leftmargin=1.5em]
\item Bernstein 分解把 $\Repinf_n$ 拆成直积 $\prod_{(L,\varsigma)}R_{(L,\varsigma)}$，
      分量的标定是「惯性等价」类 $(L,\varsigma)$，即超尖点支撑模去非分歧扭转。
\item 惯性 Langlands 对应（引理 4.4）：$\pi$ 与 $\pi'$ 在同一 Bernstein 分量
      $\Leftrightarrow$ $\rec_n(\pi)|_{I_K}\cong\rec_n(\pi')|_{I_K}$——
      即 Bernstein 中心由 Weil--Deligne 惯性类型标定。
\item 一般表示更进一步（引理 4.5）：Galois 惯性类型相同 $\Rightarrow$
      $\pi|_{\mathcal{K}}\cong\pi'|_{\mathcal{K}}$；「一般」不可去掉。
\item 类型（type）、$\mathcal{K}$-型与极小 $\mathcal{K}$-型是把 Bernstein 分量
      在 $\GL_n(\mathcal{O}_K)$ 一侧「实现」出来的工具，与 Bushnell--Kutzko 型的
      区别在于是否允许一般紧开子群。
\end{itemize}
\end{keytakeaways}
